\section{Trabalhos Relacionados}

Schwabe et al.\ \cite{schwabe2020unintended} estudaram 34 estações do sul da Califórnia (2013--2017) e relatam que a conservação de água durante a seca reduziu o volume de esgoto e aumentou sua salinidade -- o mesmo mecanismo usado aqui para interpretar a tendência de TDS da LAGWRP. Wolfand et al.\ \cite{wolfand2022dilution} modelam a bacia do rio Los Angeles, onde a LAGWRP descarrega, e relatam que o TDS aumenta com o reúso de água mesmo quando a carga total de poluentes diminui. O texto completo de nenhum dos dois foi acessível nesta sessão; a citação se apoia em resumos institucionais que confirmam a direção do efeito, não a magnitude -- nenhum número desses dois estudos é reproduzido aqui. Porse et al.\ \cite{porse2023adapting} (texto completo obtido, acesso aberto) relatam, para 45 estações do sul da Califórnia, queda de 6--10\% na vazão de efluente e aumento de 5--12\% no TDS ligados à conservação da seca de 2014--2016 -- a cifra estadual mais diretamente comparável e verificada disponível. Antweiler e Taylor \cite{antweiler2008evaluation} orientam o tratamento metodológico dos valores não detectados de BOD (Seção de Metodologia).

O estudo de maior influência sobre o enquadramento final deste trabalho é o relatório técnico da Southern California Salinity Coalition \cite{scsc2018tds} (26 estações do sul da Califórnia, texto completo lido nesta sessão): via decomposição de importância relativa (LMG), os autores atribuem a maior parte da variabilidade do TDS de influente (~88\%, contra ~12\% do consumo per capita) ao TDS da própria água de abastecimento, que varia com ciclos de seca -- não à conservação local. Esse achado reenquadra o mecanismo de Schwabe et al.\ como real, porém secundário, e motivou diretamente o teste do PDSI como variável explicativa neste trabalho (Seção de Resultados).
